% ======================================================================
% Every question type AMC supports, with its scoring (SAMPLE questions)
% Used by the sitting ../types/. See guide/08-question-types-and-scoring.md.
%
% Groups:  types_single  types_multi  types_numeric  types_open  types_shared
% ======================================================================

% ----------------------------------------------------------------------
% 1. ONE correct answer: environment "question"
%    Default score: 1 if right, 0 if wrong, blank or two boxes filled.
% ----------------------------------------------------------------------

\element{types_single}{
\begin{question}{TYPE001}
Which unit is used to measure the size of a file?
\begin{choices}
  \correctchoice{Megabyte}
  \wrongchoice{Megahertz}
  \wrongchoice{Megapixel}
  \wrongchoice{Megawatt}
\end{choices}
\end{question}
}

% Worth 2 points instead of 1.
\element{types_single}{
\begin{question}{TYPE002}
\scoring{b=2}
A table has 12 rows and 5 columns. How many cells does it have? (2 points)
\begin{choices}
  \wrongchoice{17}
  \wrongchoice{50}
  \correctchoice{60}
  \wrongchoice{125}
\end{choices}
\end{question}
}

% Wrong answers lose a quarter point; a blank answer scores 0.
\element{types_single}{
\begin{question}{TYPE003}
\scoring{b=1,m=-0.25,v=0}
Which of these is a primary key most likely to be? (a wrong answer loses 0.25)
\begin{choices}
  \correctchoice{A student ID}
  \wrongchoice{A first name}
  \wrongchoice{A birth date}
  \wrongchoice{A home town}
\end{choices}
\end{question}
}

% Choices printed in the order written, not shuffled ("o" = ordered),
% because the last one refers to the others.
\element{types_single}{
\begin{question}{TYPE004}
Which of these can be stored in a spreadsheet cell?
\begin{choices}[o]
  \wrongchoice{A number}
  \wrongchoice{A date}
  \wrongchoice{A formula}
  \correctchoice{All of the above}
\end{choices}
\end{question}
}

% Partial marks for a nearly right choice: \scoring{0.5} right after a
% choice is the score when that box is the one filled.
\element{types_single}{
\begin{question}{TYPE005}
Which chart suits monthly sales over two years best? (best answer 1, acceptable answer 0.5)
\begin{choices}
  \correctchoice{A line chart}
  \wrongchoice{A column chart}\scoring{0.5}
  \wrongchoice{A pie chart}
  \wrongchoice{A scatter plot}
\end{choices}
\end{question}
}

% ----------------------------------------------------------------------
% 2. MORE THAN ONE correct answer: environment "questionmult"
%    Students fill every box they think is right. The layout's
%    "completemulti" option adds a last box, "None of these answers are
%    correct", which is the right answer if no choice is correct.
% ----------------------------------------------------------------------

% Default: one point per box handled correctly (filled when right, left
% empty when wrong). 4 choices = up to 4 points.
\element{types_multi}{
\begin{questionmult}{TYPE101}
Which of these are file formats for images? Select all that apply.
\begin{choices}
  \correctchoice{PNG}
  \correctchoice{JPEG}
  \wrongchoice{CSV}
  \wrongchoice{MP3}
\end{choices}
\end{questionmult}
}

% All or nothing: 1 point only if every box is exactly right, otherwise 0.
\element{types_multi}{
\begin{questionmult}{TYPE102}
\scoring{mz=1}
Which of these are measures of central tendency? Select all that apply. (1 point, all or nothing)
\begin{choices}
  \correctchoice{Mean}
  \correctchoice{Median}
  \correctchoice{Mode}
  \wrongchoice{Range}
\end{choices}
\end{questionmult}
}

% 2 points, minus 1 for each box wrong, never below 0.
\element{types_multi}{
\begin{questionmult}{TYPE103}
\scoring{haut=2}
Which of these are examples of personal data? Select all that apply. (2 points, 1 off per mistake)
\begin{choices}
  \correctchoice{A phone number}
  \correctchoice{An email address}
  \wrongchoice{The population of a city}
  \wrongchoice{The average class size}
\end{choices}
\end{questionmult}
}

% Count only the right answers that were filled: +1 per correct box filled,
% -1 per wrong box filled, a box left empty counts 0, never below 0.
\element{types_multi}{
\begin{questionmult}{TYPE104}
\scoring{formula=NBC-NMC,p=0}
Which of these are cloud storage services? Select all that apply. (1 point per right answer, 1 off per wrong one)
\begin{choices}
  \correctchoice{Google Drive}
  \correctchoice{OneDrive}
  \correctchoice{Dropbox}
  \wrongchoice{Microsoft Excel}
\end{choices}
\end{questionmult}
}

% ----------------------------------------------------------------------
% 3. NUMBER answer: the student fills one digit per column.
%    \AMCnumericChoices{answer}{options}. Here: 2 digits, no sign.
% ----------------------------------------------------------------------

\element{types_numeric}{
\begin{questionmultx}{TYPE201}
Five students score 6, 7, 8, 9 and 10. What is the total of their scores?
\AMCnumericChoices{40}{digits=2,sign=false,scoreexact=1}
\end{questionmultx}
}

% A decimal answer: 1 digit before the point and 1 after.
% Close answers can earn part marks: 1 point exact, 0.5 if within 0.1.
\element{types_numeric}{
\begin{questionmultx}{TYPE202}
The mean of 2, 3 and 4.4 is what number? Give one decimal place.
\AMCnumericChoices{3.1}{digits=2,decimals=1,sign=false,scoreexact=1,scoreapprox=0.5,approx=1}
\end{questionmultx}
}

% ----------------------------------------------------------------------
% 4. WRITTEN answer, marked by hand: \AMCOpen.
%    The student writes in the lines on the answer sheet; the marker
%    fills one of the small boxes (0, 1, 2) next to it, before scanning.
% ----------------------------------------------------------------------

\element{types_open}{
\begin{question}{TYPE301}
In one sentence, explain why a backup copy should be kept in a different place from the original. (2 points, marked by the lecturer)
\AMCOpen{lines=3}{\wrongchoice[0]{0}\scoring{0}\wrongchoice[1]{1}\scoring{1}\correctchoice[2]{2}\scoring{2}}
\end{question}
}

% ----------------------------------------------------------------------
% 5. SHARED STEM: one table read by several questions.
%    The table and every question that reads it sit in ONE \element, so
%    they always print together. Draw this group whole (\DrawAll).
% ----------------------------------------------------------------------

\element{types_shared}{
\noindent\textbf{Use this table for the next two questions.}\par\smallskip
\begin{center}
\begin{tabular}{lrr}
  Shop & Mon & Tue \\ \hline
  A & 10 & 14 \\
  B & 12 & 8 \\
\end{tabular}
\end{center}

\begin{question}{TYPE401}
Which shop sold more on Tuesday?
\begin{choices}
  \correctchoice{Shop A}
  \wrongchoice{Shop B}
  \wrongchoice{They sold the same}
  \wrongchoice{It cannot be told}
\end{choices}
\end{question}

\begin{question}{TYPE402}
How many items did shop B sell over the two days?
\begin{choices}
  \wrongchoice{8}
  \wrongchoice{12}
  \correctchoice{20}
  \wrongchoice{24}
\end{choices}
\end{question}
}
